Fix an arbitrary model \(M=(G_{\mathrm w},\mathcal A_{\mathrm w},\theta)\in\mathfrak M_{\mathrm w,ZX}\), and write \(p_{\mathrm w}\) for its observed marginal law on \((Z,X,Y)\).  The parent sets in \(G_{\mathrm w}\) are
\[
\operatorname{pa}_{G_{\mathrm w}}(Z)=\varnothing,
\qquad
\operatorname{pa}_{G_{\mathrm w}}(A)=\{Z\},
\qquad
\operatorname{pa}_{G_{\mathrm w}}(X)=\{A\},
\qquad
\operatorname{pa}_{G_{\mathrm w}}(Y)=\{A,X\}.
\]
By Assumption \(\mathrm{ass:qualitative\text{-}functionality}\), for each \(z\in\{0,1\}\) the row \(\theta_A(\cdot\mid z)\) is a point mass.  Hence there is a unique map \(f:\{0,1\}\to\{0,1\}\) such that
\[
\theta_A(a\mid z)=\mathbf 1\{a=f(z)\}
\qquad(a,z\in\{0,1\}).
\]

Assumption \(\mathrm{ass:causal\text{-}factorization}\), specialized to the displayed parent sets, gives, for every \(z,a,x,y\in\{0,1\}\),
\[
\begin{aligned}
P_M(Z=z,A=a,X=x,Y=y)
&=\theta_Z(z)\theta_A(a\mid z)\theta_X(x\mid a)\theta_Y(y\mid a,x)\\
&=\theta_Z(z)\mathbf 1\{a=f(z)\}\theta_X(x\mid a)\theta_Y(y\mid a,x).
\end{aligned}
\]
Marginalizing the hidden coordinate \(A\) therefore yields
\[
\begin{aligned}
p_{\mathrm w}(z,x,y)
&=\sum_{a\in\{0,1\}}P_M(Z=z,A=a,X=x,Y=y)\\
&=\theta_Z(z)\theta_X(x\mid f(z))\theta_Y(y\mid f(z),x).
\end{aligned}
\]
Because every conditional-probability row is normalized, summing this identity over \(y'\in\{0,1\}\) gives
\[
\begin{aligned}
p_{\mathrm w}(z,x)
&=\sum_{y'\in\{0,1\}}p_{\mathrm w}(z,x,y')\\
&=\theta_Z(z)\theta_X(x\mid f(z))
  \sum_{y'\in\{0,1\}}\theta_Y(y'\mid f(z),x)\\
&=\theta_Z(z)\theta_X(x\mid f(z)).
\end{aligned}
\]
The support premise in the theorem states that this last quantity is strictly positive for every \((z,x)\).  The same guard also follows from Assumption \(\mathrm{ass:joint\text{-}treatment\text{-}positivity}\) after specializing \(T\) to \(T_{ZX}=\{Z,X\}\).  Thus the observable functional
\[
\Phi(p_{\mathrm w})(z,x,y)
:=
\frac{p_{\mathrm w}(z,x,y)}{p_{\mathrm w}(z,x)}
\]
is well defined on every model under consideration.  Dividing the preceding two factorizations, with the strict-support guard just established, gives
\[
\Phi(p_{\mathrm w})(z,x,y)
=\theta_Y(y\mid f(z),x)
=p_{\mathrm w}(Y=y\mid Z=z,X=x).
\]
This defines only an observed-law functional; the next calculation proves that it equals the causal target.

Under \(do(Z=z,X=x)\), Assumption \(\mathrm{ass:truncated\text{-}intervention}\) replaces the rows at \(Z\) and \(X\) by point masses while leaving the mechanisms at \(A\) and \(Y\) unchanged.  Marginalizing the non-outcome coordinate \(A\) in that truncated product gives
\[
\begin{aligned}
Q_{ZX,M}(z,x,y)
&=\sum_{a\in\{0,1\}}\theta_A(a\mid z)\theta_Y(y\mid a,x)\\
&=\sum_{a\in\{0,1\}}\mathbf 1\{a=f(z)\}\theta_Y(y\mid a,x)\\
&=\theta_Y(y\mid f(z),x).
\end{aligned}
\]
Combining the last two displays proves, for every \(z,x,y\in\{0,1\}\),
\[
Q_{ZX,M}(z,x,y)
=\Phi(p_{\mathrm w})(z,x,y)
=p_{\mathrm w}(Y=y\mid Z=z,X=x).
\]

Finally, let \(M,M'\in\mathfrak M_{\mathrm w,ZX}\) induce the same observed law \(p_{\mathrm w}\).  Applying the established equality to each model gives
\[
Q_{ZX,M}(z,x,y)=\Phi(p_{\mathrm w})(z,x,y)=Q_{ZX,M'}(z,x,y)
\]
for every \(z,x,y\).  Hence the entire joint-intervention distribution is uniquely determined by the observed law, which is precisely graphwide identification in this class.  Every division in the recovery map is licensed by the explicitly proved guard \(p_{\mathrm w}(z,x)>0\).